\documentclass[10pt,conference]{IEEEtran}

\usepackage[utf8]{inputenc}
\usepackage{cite}
\usepackage{amsmath,amssymb,amsfonts}
\usepackage{algorithmic}
\usepackage{graphicx}
\usepackage{textcomp}
\usepackage{xcolor}
\usepackage{booktabs}
\usepackage{multirow}
\usepackage{array}
\usepackage{url}
\usepackage{hyperref}
\hypersetup{colorlinks=true,linkcolor=blue,citecolor=blue,urlcolor=blue}

\def\BibTeX{{\rm B\kern-.05em{\sc i\kern-.025em b}\kern-.08em
    T\kern-.1667em\lower.7ex\hbox{E}\kern-.125emX}}

\begin{document}

\title{Event Cameras Meet Gaussian Splatting: A Survey of Neuromorphic 3D and 4D Reconstruction}

\author{
\IEEEauthorblockN{Arthur Li Rui}
\IEEEauthorblockA{\textit{ERF-GS Research Project} \\
\textit{Independent Survey}\\
arthurlirui@example.com}
}

\maketitle

\begin{abstract}
The intersection of event cameras and Gaussian Splatting (GS) has emerged as one of the most active frontiers in novel-view synthesis and dynamic scene reconstruction. Event cameras asynchronously report per-pixel log-luminance changes with microsecond latency, high dynamic range, and negligible motion blur, complementing the limitations of frame-based RGB sensors. Gaussian Splatting, with its explicit, real-time, differentiable rasterization of anisotropic primitives, offers an ideal renderer onto which event-based supervision can be grafted. This survey provides a comprehensive review of the rapidly growing literature on \emph{Event-based Gaussian Splatting} (Event-GS), spanning 2023--2026, with particular emphasis on recent CVPR, ICCV, ECCV, NeurIPS, and ICML work. We organize $24+$ methods into six taxonomic categories---static novel-view synthesis, motion deblurring, low-light/HDR reconstruction, dynamic/4D reconstruction, SLAM/tracking, and dataset/simulation---and analyze the core technical contribution, event-generation model, deformation strategy, and supervision formulation of each. We further extract cross-cutting technical dimensions (event-to-image vs.\ event-to-flow supervision, continuous trajectory parameterization, deformation fields, physical priors) and identify open problems including occlusion handling, static-prior bootstrapping, per-event point-process losses, and differentiable event-camera physics. This survey aims to serve as a reference for researchers entering the field and as a roadmap for future work.
\end{abstract}

\begin{IEEEkeywords}
event cameras, Gaussian Splatting, novel view synthesis, dynamic scene reconstruction, deblurring, neuromorphic vision, 4D reconstruction, survey
\end{IEEEkeywords}

%==================================================================
\section{Introduction}
\label{sec:introduction}

Novel-view synthesis (NVS) and 3D/4D scene reconstruction have been transformed by two complementary technologies. On the representation side, 3D Gaussian Splatting (3DGS)~\cite{kerbl2023gaussian} replaced the implicit neural radiance fields of NeRF~\cite{mildenhall2021nerf} with an explicit set of anisotropic Gaussian primitives that can be rasterized in real time and optimized via differentiable rendering. On the sensing side, event cameras~\cite{gallego2020event}---bio-inspired sensors that asynchronously emit per-pixel log-luminance change events---provide microsecond temporal resolution, high dynamic range (HDR), and near-immunity to motion blur, precisely the capabilities that conventional frame cameras lack.

The combination is natural. RGB frames supply dense absolute appearance but blur and alias under fast motion; events supply sparse, relative, high-temporal-resolution motion cues but no absolute color. Recent work has demonstrated that fusing the two modalities through the Gaussian Splatting renderer yields reconstructions that neither modality can achieve alone: sharp novel views from motion-blurred inputs~\cite{yu2025evagaussians,lee2025dietgs}, HDR radiance fields in the dark~\cite{wu2025darkevgs,chen2026evhdrgs}, and dynamic 4D Gaussians for fast-moving scenes~\cite{dai2026fasteventdgs,bai2026erfgs,xu2025eventboosted}.

The pace of progress has been remarkable. The first event-based Gaussian Splatting methods appeared in 2024~\cite{han2024event3dgs,deguchi2024e2gs,wang2024evggs}; by 2025--2026 the subfield spans dedicated tracks at CVPR, ICCV, ECCV, NeurIPS, and ICML, with methods targeting deblurring (EvaGaussians~\cite{yu2025evagaussians}, DiET-GS~\cite{lee2025dietgs}), free-trajectory reconstruction (EF-3DGS~\cite{liao2025ef3dgs}, EvTrajGS~\cite{chen2026evtrajgs}), fast dynamic 4D (FastEventDGS~\cite{dai2026fasteventdgs}, ERF-GS~\cite{bai2026erfgs}, Event-Boosted~\cite{xu2025eventboosted}), and SLAM/tracking (GS-EVT~\cite{liu2025gsevt}, MotionGS-SLAM~\cite{hu2026motiongs}). The diversity of formulations---event double integral priors, contrast maximization, linear-flow losses, continuous B-spline trajectories, event-driven densification---makes a consolidated survey timely.

\textbf{Scope and contributions.} This survey reviews the literature on event-camera-driven Gaussian Splatting through mid-2026. Our contributions are:
\begin{itemize}
    \item A six-category taxonomy (\S\ref{sec:taxonomy}) organizing $24+$ methods by the sub-problem they solve.
    \item Per-paper analysis (\S\ref{sec:static}--\S\ref{sec:datasets}) of the core technical contribution, event-generation model, and supervision form of each method, with emphasis on CVPR/ICCV/ECCV work.
    \item A cross-cutting technical synthesis (\S\ref{sec:crosscutting}) extracting the recurring design choices (event-to-image vs.\ event-to-flow, trajectory parameterization, deformation fields, physical priors).
    \item An open-problems discussion (\S\ref{sec:discussion}) identifying occlusion handling, static-prior bootstrapping, per-event point-process losses, and differentiable event-camera physics as the most promising future directions.
\end{itemize}

We deliberately complement event-camera-only NeRF works (EventNeRF~\cite{rudnev2023eventnerf}, E-NeRF~\cite{klenk2023enerf}, Ev-NeRF~\cite{hwang2023evnerf}) only as background, and focus the bulk of the survey on the Gaussian-Splatting family, which now dominates the field in rendering speed and reconstruction fidelity.

\section{Background}
\label{sec:background}

\subsection{Event Cameras}
\label{sec:bg:event}

An event camera (Dynamic Vision Sensor, DVS) emits an event $e_k=(x_k,y_k,t_k,p_k)$ whenever the log-luminance $L$ at pixel $(x_k,y_k)$ changes by a contrast threshold $C$:
\begin{equation}
    L(x,y,t_k) - L(x,y,t_k-\Delta t) = p_k\, C, \quad p_k\in\{-1,+1\}.
    \label{eq:event}
\end{equation}
Events are sparse, asynchronous, and carry only relative brightness. Their strengths---microsecond temporal resolution ($\sim\!\mu$s), HDR ($>120$\,dB), no motion blur---are dual to their weaknesses: no absolute intensity, noisy in low-contrast regions, and difficult to aggregate into a dense representation. The contrast threshold $C$ is in practice pixel-, time-, and scene-dependent, a non-ideality that motivates the differentiable simulator line of work discussed in \S\ref{sec:discussion}.

\textbf{Event generation model.} The canonical mapping from a pair of images $I_1,I_2$ to a predicted event stream is the \emph{linlog} model:
\begin{equation}
    \hat{e} = \frac{\mathrm{linlog}(\mathrm{luma}(I_2)) - \mathrm{linlog}(\mathrm{luma}(I_1))}{C},
    \label{eq:linlog}
\end{equation}
where $\mathrm{luma}=0.299R+0.587G+0.114B$ and $\mathrm{linlog}$ is linear below a threshold and logarithmic above, mimicking the sensor response~\cite{bai2026erfgs}. Aggregated forms include the event voxel (positive/negative polarity counts binned into $T$ time slices) and the event double integral (EDI)~\cite{pan2019edi}, which relates a blurry image to a sequence of sharp images via integration of the event stream.

\subsection{Gaussian Splatting}
\label{sec:bg:gs}

3DGS~\cite{kerbl2023gaussian} represents a scene as a set of anisotropic Gaussians $\{\mu_i,\Sigma_i,\alpha_i,c_i\}$ (mean, covariance, opacity, color via spherical harmonics). Rendering is an $\alpha$-blended differentiable rasterization:
\begin{equation}
    C(\mathbf{u}) = \sum_{i\in\mathcal{N}} c_i\,\alpha_i\,T_i, \quad T_i=\prod_{j<i}(1-\alpha_j),
    \label{eq:splat}
\end{equation}
where $\alpha_i=o_i\cdot\exp(-\tfrac12(\mathbf{u}-\mu_i')^\top\Sigma_i^{-1}(\mathbf{u}-\mu_i'))$ after projecting the 3D Gaussian to the screen. Optimization adapts the per-Gaussian parameters via gradient-driven densification (clone/split), pruning, and opacity reset.

\textbf{Deformable / 4D GS.} Dynamic extensions add a deformation field $D(\mathbf{x},t)$ that predicts per-Gaussian deltas $\Delta\mu,\Delta\Sigma,\Delta\alpha,\Delta c$ at query time $t$. Representative backbones include deformation MLPs conditioned on HexPlane~\cite{cao2023hexplane} grid features (4DGS~\cite{wu2024ldgswu}), per-Gaussian embeddings~\cite{bae2024pergaussian}, and 4D Gaussian primitives~\cite{yang2024realtime}. These backbones are the substrate onto which event-camera supervision is grafted in the dynamic-reconstruction methods of \S\ref{sec:dynamic}.

\subsection{Why Events + GS?}
\label{sec:bg:why}

The pairing is synergistic along three axes:
\begin{itemize}
    \item \emph{Speed:} 3DGS renders in real time ($>100$\,FPS), enabling the many-frame event-supervision loops that NeRF could not afford.
    \item \emph{Differentiability:} The differentiable rasterizer lets event losses backpropagate into Gaussian geometry and appearance directly, without the surrogate losses required by NeRF.
    \item \emph{Explicit densification:} 3DGS's clone/split machinery offers a natural insertion point for \emph{event-driven} densification---back-projecting event pixels to spawn new primitives---which has no analog in implicit representations.
\end{itemize}
These three properties jointly explain why the Event-GS literature has grown from a single 2024 method to a multi-track subfield in under three years.

\section{Taxonomy}
\label{sec:taxonomy}

We organize the Event-GS literature into six categories by the sub-problem addressed. A method may appear in more than one category; we place each in its primary category and cross-reference.

\begin{enumerate}
    \item[\textbf{A.}] \textbf{Static novel-view synthesis} (\S\ref{sec:static}): event-only or event-RGB reconstruction of static scenes, often with fast camera motion or unknown pose. Key works: EventSplat~\cite{yura2025eventsplat}, Event-3DGS~\cite{han2024event3dgs}, EvGGS~\cite{wang2024evggs}, EV-GS~\cite{wu2024evgs}, EF-3DGS~\cite{liao2025ef3dgs}, EvTrajGS~\cite{chen2026evtrajgs}, IncEventGS~\cite{huang2025inceventgs}, E2EGS~\cite{kim2026e2egs}, GS-EVT~\cite{liu2025gsevt}.
    \item[\textbf{B.}] \textbf{Motion deblurring} (\S\ref{sec:deblurring}): recovering sharp 3DGS from motion-blurred RGB frames aided by events. Key works: E2GS~\cite{deguchi2024e2gs}, EvaGaussians~\cite{yu2025evagaussians}, DiET-GS~\cite{lee2025dietgs}, EBAD-Gaussian~\cite{deng2025ebad}, EvFlow-GS~\cite{an2026evflow}, JADE-GS~\cite{fu2026jadegs}, E2D-GS~\cite{lin2025e2dgs}, EaDeblur-GS~\cite{weng2024eadeblur}, Ev3DGS~\cite{huang2024ev3dgs}, DeblurSplat~\cite{li2025deblursplat}, AsyncEvGS~\cite{dai2026asyncevgs}, BAD-Gaussians~\cite{zhao2024badgaussians}.
    \item[\textbf{C.}] \textbf{Low-light / HDR} (\S\ref{sec:lowlight}): reconstructing radiance fields under extreme illumination. Key works: Dark-EvGS~\cite{wu2025darkevgs}, EvHDR-GS~\cite{chen2026evhdrgs}.
    \item[\textbf{D.}] \textbf{Dynamic / 4D reconstruction} (\S\ref{sec:dynamic}): deformable Gaussians for fast-moving scenes. Key works: FastEventDGS~\cite{dai2026fasteventdgs}, ERF-GS~\cite{bai2026erfgs}, Event-Boosted~\cite{xu2025eventboosted}, PEGS~\cite{xu2025pegs}.
    \item[\textbf{E.}] \textbf{SLAM / tracking} (\S\ref{sec:slam}): event-driven localization and mapping with GS maps. Key works: MotionGS-SLAM~\cite{hu2026motiongs}, GS-EVT~\cite{liu2025gsevt}.
    \item[\textbf{F.}] \textbf{Datasets / simulation} (\S\ref{sec:datasets}): using GS to generate event data, or benchmark datasets. Key works: GS2E~\cite{li2025gs2e}.
\end{enumerate}

Table~\ref{tab:summary} summarizes the methods, venues, input modalities, and key innovations.

\begin{table*}[t]
\centering
\caption{Summary of Event-based Gaussian Splatting methods (2023--2026). ``In'' = input modality (E=events, R=RGB, B=blurry RGB, D=depth). ``Pose'' = whether known poses are required (\checkmark) or estimated (Pose-Free).}
\label{tab:summary}
\renewcommand{\arraystretch}{1.15}
\small
\begin{tabular}{@{}p{0.5cm}p{3.4cm}p{1.2cm}p{0.7cm}p{0.7cm}p{8.5cm}@{}}
\toprule
\# & Method & Venue & In & Pose & Core innovation \\
\midrule
\multicolumn{6}{l}{\textbf{A. Static novel-view synthesis}} \\
A1 & Event-3DGS~\cite{han2024event3dgs} & NeurIPS'24 & E & \checkmark & First event-only 3DGS; view-volume loss for stability. \\
A2 & EvGGS~\cite{wang2024evggs} & ICML'24 & E & \checkmark & First generalizable feed-forward event-GS; collaborative learning. \\
A3 & EV-GS~\cite{wu2024evgs} & MVA'24 & E & \checkmark & CNI-informed monocular event-GS. \\
A4 & EventSplat~\cite{yura2025eventsplat} & CVPR'25 & E & \checkmark & E2VID init + spline-interpolated poses; real-time. \\
A5 & IncEventGS~\cite{huang2025inceventgs} & CVPR'25 & E & Free & Incremental pose-free event-GS; SLAM-style tracking+mapping. \\
A6 & EF-3DGS~\cite{liao2025ef3dgs} & NeurIPS'25 & E,R & Free & Free-trajectory; CMax pose calibration + photometric BA. \\
A7 & GS-EVT~\cite{liu2025gsevt} & ICRA'25 & E,R & \checkmark & Cross-modal: RGB-built GS map, event-camera tracking. \\
A8 & EvTrajGS~\cite{chen2026evtrajgs} & arXiv'26 & E & Free & Continuous-time trajectory + temporally-coupled pose; 7$\times$ faster than IncEventGS. \\
A9 & E2EGS~\cite{kim2026e2egs} & arXiv'26 & E & Free & Event-to-edge pose-free reconstruction. \\
\midrule
\multicolumn{6}{l}{\textbf{B. Motion deblurring}} \\
B1 & E2GS~\cite{deguchi2024e2gs} & ICIP'24 & B,E & \checkmark & First event+GS; EDI bins; 140 FPS, 60$\times$ faster than E$^2$NeRF. \\
B2 & Ev3DGS~\cite{huang2024ev3dgs} & APSIPA'24 & B,E & \checkmark & Event+RGB 3DGS deblurring. \\
B3 & EaDeblur-GS~\cite{weng2024eadeblur} & arXiv'24 & B,E & \checkmark & Event-assisted real-time deblur 3DGS. \\
B4 & BAD-Gaussians~\cite{zhao2024badgaussians} & ECCV'24 & B & \checkmark & Bundle-adjusted deblur GS (no events; baseline). \\
B5 & EvaGaussians~\cite{yu2025evagaussians} & ICCV'25 & B,E & \checkmark & Event streams into init+optimization; two new datasets. \\
B6 & DiET-GS~\cite{lee2025dietgs} & CVPR'25 & B,E & \checkmark & EDI + diffusion prior (RSD); two-stage deblurring. \\
B7 & EBAD-Gaussian~\cite{deng2025ebad} & arXiv'25 & B,E & Est. & Event-driven bundle-adjusted deblur; EDI prior. \\
B8 & EvFlow-GS~\cite{an2026evflow} & ICME'26 & B,E & Est. & Learnable double integral + optical flow; B\'ezier SE(3) trajectory; critiques linlog. \\
B9 & DeblurSplat~\cite{li2025deblursplat} & arXiv'25 & B,E & Free & First SfM-free deblur event-GS. \\
B10 & JADE-GS~\cite{fu2026jadegs} & arXiv'26 & B,E & \checkmark & Spatial prior router fuses analytic+learned deblur evidence. \\
B11 & E2D-GS~\cite{lin2025e2dgs} & ESWA'25 & B,E & \checkmark & Event-driven blur formation + differential consistency. \\
B12 & AsyncEvGS~\cite{dai2026asyncevgs} & arXiv'26 & B,E & \checkmark & Asynchronous RGB-Event dual-camera system + dataset. \\
\midrule
\multicolumn{6}{l}{\textbf{C. Low-light / HDR}} \\
C1 & Dark-EvGS~\cite{wu2025darkevgs} & arXiv'25 & B,E & \checkmark & Triplet supervision + color tone matching; first dark-event-radiance dataset. \\
C2 & EvHDR-GS~\cite{chen2026evhdrgs} & AAAI'26 & E,R & \checkmark & HDR 3DGS; learnable HDR-to-LDR via events. \\
\midrule
\multicolumn{6}{l}{\textbf{D. Dynamic / 4D reconstruction}} \\
D1 & Event-Boosted~\cite{xu2025eventboosted} & ICCV'25 & E,R & \checkmark & GS-Threshold joint modeling + dynamic-static decomposition. \\
D2 & FastEventDGS~\cite{dai2026fasteventdgs} & CVPR'26 & E & \checkmark & First single-event-camera deformable 4DGS; B-spline + dual event model + VGGT depth. \\
D3 & ERF-GS~\cite{bai2026erfgs} & CVMJ'26 & E,R & \checkmark & Disjoint event-RGB viewpoints; linlog+STE loss; event-driven densification. \\
D4 & PEGS~\cite{xu2025pegs} & arXiv'25 & E,R & \checkmark & Physics prior (acceleration consistency) + Kalman; large-scale rigid motion. \\
\midrule
\multicolumn{6}{l}{\textbf{E. SLAM / tracking}} \\
E1 & MotionGS-SLAM~\cite{hu2026motiongs} & arXiv'26 & E,R & Free & Forward blur modeling; event-modulated anisotropic kernels; 4D hash grid. \\
\midrule
\multicolumn{6}{l}{\textbf{F. Datasets / simulation}} \\
F1 & GS2E~\cite{li2025gs2e} & NeurIPS'25 & R & \checkmark & GS$\to$events simulator; 1150+ scenes; adaptive threshold. \\
\bottomrule
\end{tabular}
\end{table*}

\section{Static Novel-View Synthesis}
\label{sec:static}

This category covers methods that reconstruct a static 3D Gaussian scene from event-camera input, typically with a fast-moving camera and possibly unknown poses. The central challenge is that events carry no absolute intensity, so photometric supervision must be synthesized indirectly.

\subsection{Event-3DGS (NeurIPS 2024)}
\label{sec:static:event3dgs}

Event-3DGS~\cite{han2024event3dgs} is the first event-based 3DGS method. It optimizes Gaussian parameters from event streams by rendering images at adjacent timestamps and enforcing consistency between the rendered log-luminance difference and the integral of events in the interval. To stabilize optimization under the sparse, noisy event signal, it introduces a \emph{view-volume consistency loss} that penalizes large spatial derivatives of the rendered event image. It demonstrates superior reconstruction quality over event-NeRF baselines on both simulated and real datasets, while inheriting 3DGS's real-time rendering.

\textbf{Core point:} establish that differentiable GS rasterization can be supervised directly by event integrals, opening the Event-GS paradigm.

\subsection{EvGGS (ICML 2024)}
\label{sec:static:evggs}

EvGGS~\cite{wang2024evggs} is the first \emph{generalizable} event-GS: rather than per-scene optimization, it reconstructs 3D Gaussians from event input in a single feed-forward pass and generalizes to unseen scenes without retraining. It uses a collaborative learning framework with multiple losses (image, feature, geometry) and an event-to-image regression sub-network that provides pseudo-RGB supervision. A distillation strategy transfers the regression network's knowledge to the GS renderer.

\textbf{Core point:} move Event-GS from per-scene optimization to feed-forward generalization, a direction RGB GS explored later with pixel-splat methods.

\subsection{EV-GS (MVA 2024)}
\label{sec:static:evgs}

EV-GS~\cite{wu2024evgs} is the first Contrast-Noise-Intensity (CNI)-informed scheme for monocular event-GS. It explicitly models the CNI relationship of the event sensor to improve the fidelity of the predicted event image from rendered Gaussians, enabling efficient novel-view synthesis from a single moving event camera.

\subsection{EventSplat (CVPR 2025)}
\label{sec:static:eventsplat}

EventSplat~\cite{yura2025eventsplat} targets static-scene NVS from a fast-moving event camera. Two design choices distinguish it: (i) initialization via an event-to-video (E2VID) model produces a denser, more accurate initial point cloud than naive event-voxel back-projection; (ii) spline interpolation across the event-camera trajectory yields high-quality per-event poses, mitigating the pose-aliasing that plagues fast-motion event capture. It achieves real-time rendering and outperforms event-NeRF by an order of magnitude in speed.

\textbf{Core point:} initialization quality and pose continuity dominate final fidelity; E2VID + splines are simple, effective priors.

\subsection{IncEventGS (CVPR 2025)}
\label{sec:static:inceventgs}

IncEventGS~\cite{huang2025inceventgs} is the first \emph{pose-free} event-GS from a single event camera. It runs an incremental tracking-and-mapping loop: events track camera motion via differentiable rendering of the current Gaussian map, while mapping updates Gaussians with the tracked poses. It outperforms NeRF-based pose-free event methods even without ground-truth poses, at the cost of SLAM-style computational overhead.

\subsection{EF-3DGS (NeurIPS 2025, Spotlight)}
\label{sec:static:ef3dgs}

EF-3DGS~\cite{liao2025ef3dgs} introduces the event camera to \emph{free-trajectory} 3DGS---reconstruction from a casually captured video with no precise poses. It combines three components: (1) an Event Generation Model that fuses events and frames for continuous inter-frame supervision; (2) Contrast Maximization (CMax) for joint pose calibration and gradient-domain 3DGS constraints; (3) photometric Bundle Adjustment with a fixed-GS separation of structure and color optimization. It reports $+3$\,dB PSNR and $-40\%$ ATE on high-speed scenes versus pose-free RGB baselines.

\textbf{Core point:} events are a natural regularizer for pose-free reconstruction; CMax gives a calibration signal that RGB-only SfM lacks under fast motion.

\subsection{GS-EVT (ICRA 2025)}
\label{sec:static:gsevt}

GS-EVT~\cite{liu2025gsevt} is a cross-modal tracking method: a 3DGS map is built from RGB frames, then an event camera tracks against it. It renders local differential images from the GS map using a first-order motion prior and reference poses, then matches them to integrated event maps via coarse-to-fine optimization. This decouples map-building (RGB, dense) from tracking (events, robust to fast motion).

\subsection{EvTrajGS (arXiv 2026)}
\label{sec:static:evtrajgs}

EvTrajGS~\cite{chen2026evtrajgs} addresses the efficiency--fidelity trade-off of pose-free event-GS. It parameterizes camera motion as a continuous-time trajectory initialized from coarse pose priors, then aggregates adjacent trajectory states into a \emph{temporally coupled pose} that promotes consistent joint pose--scene updates---avoiding the expensive incremental SLAM loop of IncEventGS. A loss-reweighted event sampling strategy emphasizes under-reconstructed temporal intervals. It achieves $+3.8$\,dB PSNR, $+0.1$ SSIM, and $>40\%$ lower ATE while running $7\times$ faster than IncEventGS.

\textbf{Core point:} continuous-time trajectory parameterization replaces SLAM-style incremental tracking with a unified joint optimization, a recurring design in later work.

\subsection{E2EGS (arXiv 2026)}
\label{sec:static:e2egs}

E2EGS~\cite{kim2026e2egs} is a pose-free framework operating solely on event streams. Its name reflects the use of event edges as the primary geometric cue: it aligns rendered edge maps to event-edge observations for joint pose--scene optimization, achieving strong reconstruction quality and trajectory accuracy and establishing a fully pose-free event-only paradigm.

\section{Motion Deblurring}
\label{sec:deblurring}

Motion deblurring is the most populated Event-GS category. The setup is canonical: a moving camera captures blurry RGB frames during prolonged exposure, while a co-located event camera records blur-free log-luminance changes. The Event Double Integral (EDI)~\cite{pan2019edi} prior relates a blurry image $B$ to a sequence of sharp images $\{I_k\}$ and the event stream:
\begin{equation}
    B = \frac{1}{N}\sum_{k=1}^{N} I_k, \quad I_k = I_0 \cdot \exp\!\Big(\int_0^{t_k} e(\tau)\,d\tau\Big).
    \label{eq:edi}
\end{equation}
Almost every method in this section is a variant of Eq.~\eqref{eq:edi} grafted onto differentiable GS rendering, differing in how the event prior is parameterized (analytic vs.\ learnable), whether poses are estimated, and what auxiliary priors (diffusion, optical flow) are added.

\subsection{E2GS (ICIP 2024)}
\label{sec:deblur:e2gs}

E2GS~\cite{deguchi2024e2gs} is the first work to incorporate event data into Gaussian Splatting. It combines blurry images with events for deblurring NVS, using the EDI model to split the exposure into $N$ event bins and estimate intermediate sharp frames. It achieves 140\,FPS rendering and is $60\times$ faster to train and $3500\times$ faster to render than the prior event-NeRF baseline E$^2$NeRF.

\textbf{Core point:} establish the event+GS deblurring paradigm and the EDI-bin supervision form.

\subsection{EvaGaussians (ICCV 2025)}
\label{sec:deblur:evagaussians}

EvaGaussians~\cite{yu2025evagaussians} integrates event streams into both the \emph{initialization} and \emph{optimization} of 3DGS. Events help construct a sharper initial point cloud and provide continuous inter-frame supervision during optimization. It contributes two new RGB+event+pose datasets and demonstrates improved fidelity and training/inference efficiency over EDI-only baselines.

\subsection{DiET-GS (CVPR 2025)}
\label{sec:deblur:dietgs}

DiET-GS~\cite{lee2025dietgs} extends the EDI prior in two directions. First, it models EDI in the \emph{brightness domain} via a learnable camera response function, rather than treating RGB channels as brightness directly---this better adapts to real sensor non-linearities and recovers finer details. Second, it injects a \emph{diffusion prior} via Renoised Score Distillation (RSD) in a two-stage pipeline: stage~1 (DiET-GS) jointly uses real events and diffusion; stage~2 (DiET-GS++) adds learnable latent residuals per Gaussian and optimizes with diffusion alone to sharpen edges.

\textbf{Core point:} generative priors (diffusion) complement event priors; the two-stage schedule prevents the diffusion prior from being diluted by real-data loss.

\subsection{EBAD-Gaussian (arXiv 2025)}
\label{sec:deblur:ebad}

EBAD-Gaussian~\cite{deng2025ebad} reconstructs sharp 3D Gaussians from severely blurred images and event streams via an event-driven Bundle Adjustment. It synthesizes multiple sharp intermediate frames within the exposure to construct a blur loss, supervises arbitrary-time brightness change via events, and constrains intermediate frames with the EDI prior. It jointly learns Gaussian parameters and recovers the intra-exposure camera trajectory.

\subsection{EvFlow-GS (ICME 2026)}
\label{sec:deblur:evflowgs}

EvFlow-GS~\cite{an2026evflow} is notable for explicitly \emph{critiquing} the linlog+fixed-threshold event-generation model on four grounds: noise amplification, constant-threshold assumption, low-intensity-region degradation, and implicit constant-velocity assumption. It replaces it with an end-to-end \emph{learnable double integral (LDI)} supervised by optical flow: events are warped by optical flow to align motion trajectories, yielding sharp supervision. Camera motion is modeled with a 9-control-point B\'ezier SE(3) trajectory.

\textbf{Core point:} a methodological inflection---the field begins to question the linlog+threshold approximation and move toward learnable, flow-based, or physics-based event models.

\subsection{JADE-GS (arXiv 2026)}
\label{sec:deblur:jadegs}

JADE-GS~\cite{fu2026jadegs} treats deblurring as \emph{evidence allocation}: two families of priors---the analytic EDI inversion and a learned frame+event restoration network---are viewed as complementary evidence sources. A lightweight Spatial Prior Router performs pixel-level allocation to fuse them into an auxiliary supervision target. The router trains without clean reference frames and is removed after optimization. JADE-GS leads LPIPS and CLIP-IQA on both synthetic and real benchmarks while training in $\sim$1\,h on a consumer GPU.

\subsection{E2D-GS (ESWA 2025)}
\label{sec:deblur:e2dgs}

E2D-GS~\cite{lin2025e2dgs} physically models motion-blur formation from event streams and optimizes the discrepancy between synthesized and observed blurry images while recovering the camera trajectory. A differential-consistency module denoises events and regularizes Gaussian parameters for robustness in real-world conditions.

\subsection{EaDeblur-GS and Ev3DGS (2024)}
\label{sec:deblur:early}

EaDeblur-GS~\cite{weng2024eadeblur} integrates event data into 3DGS for real-time sharp reconstruction from blurred inputs, achieving comparable quality to state-of-the-art with real-time rendering. Ev3DGS~\cite{huang2024ev3dgs} is a concurrent APSIPA'24 framework combining event and RGB cameras with 3DGS, improving rendering performance and reducing training time.

\subsection{DeblurSplat (arXiv 2025) and AsyncEvGS (arXiv 2026)}
\label{sec:deblur:free}

DeblurSplat~\cite{li2025deblursplat} is the first SfM-free deblurring event-GS: it removes the Structure-from-Motion prerequisite, achieving high-fidelity novel views with significant rendering efficiency. AsyncEvGS~\cite{dai2026asyncevgs} introduces a flexible, high-resolution \emph{asynchronous} RGB-Event dual-camera system and contributes the AsyncEv-Deblur dataset captured with this hardware, addressing the synchronization bottleneck of rigidly coupled RGB-event rigs.

\subsection{BAD-Gaussians (ECCV 2024, non-event baseline)}
\label{sec:deblur:bad}

BAD-Gaussians~\cite{zhao2024badgaussians} is a non-event bundle-adjusted deblur 3DGS that handles severe motion blur with inaccurate poses. We include it as the canonical \emph{non-event} baseline against which the event-assisted methods above are benchmarked; it establishes the explicit-Gaussian deblurring formulation that event priors later augment.

\section{Low-Light and HDR Reconstruction}
\label{sec:lowlight}

Event cameras are uniquely suited to extreme illumination: their HDR ($>120$\,dB) and logarithmic response preserve information in both dark and bright regions where RGB frames saturate or are photon-starved. Two methods pioneeringly exploit this for radiance-field reconstruction.

\subsection{Dark-EvGS (arXiv 2025)}
\label{sec:lowlight:darkevgs}

Dark-EvGS~\cite{wu2025darkevgs} is the first event-guided dark-light 3DGS bright-frame synthesis. Its core is \emph{triplet-level supervision}: three complementary losses constrain (i) overall scene structure, (ii) fine-grained details, and (iii) sharp rendering. A Color Tone Matching Block ensures color consistency between the synthesized bright frames and the dark RGB observations, addressing the color-shift artifact that naive event-to-image supervision produces. It contributes the first real-world dark-light event-radiance-field dataset.

\textbf{Core point:} in low light, the event signal dominates the RGB signal; the challenge shifts from deblurring to color and detail recovery, requiring explicit tone-matching.

\subsection{EvHDR-GS (AAAI 2026)}
\label{sec:lowlight:evhdrgs}

EvHDR-GS~\cite{chen2026evhdrgs} reframes HDR video reconstruction as HDR \emph{scene} reconstruction: a 3DGS representation guarantees 3D-consistent temporal brightness. HDR 3D Gaussians simultaneously represent HDR and LDR color, and a learnable HDR-to-LDR tone-mapping is self-supervised by the event stream, removing the dataset bias of fixed tone-mapping operators. The event stream acts as the bridge between HDR radiance and LDR observations.

\textbf{Core point:} the event stream's logarithmic response is a natural differentiable tone-mapping supervisor; coupling it with explicit 3DGS consistency yields temporally stable HDR NVS.

\section{Dynamic and 4D Reconstruction}
\label{sec:dynamic}

Dynamic-scene Event-GS is the most methodologically active category. The common substrate is a Deformable 3DGS (deformation MLP + HexPlane grid, or 4D primitives) that predicts per-Gaussian deltas at query time $t$. Events provide the high-temporal-resolution motion supervision that RGB cannot, but raise three challenges: (i) how to query arbitrary timestamps given sparse poses; (ii) how to constrain photometry without absolute intensity; (iii) how to correct depth under monocular sparse views.

\subsection{Event-Boosted Deformable 3D Gaussians (ICCV 2025)}
\label{sec:dynamic:boosted}

Event-Boosted~\cite{xu2025eventboosted} is the first event-camera + deformable-3DGS combination for dynamic scenes. It makes two contributions. (1) \emph{GS-Threshold Joint Modeling (GTJM)}: observing that event threshold modeling is critical to reconstruction quality, it jointly optimizes the Gaussian field and the event threshold, with each informing the other---eliminating the ``purple haze'' artifact caused by a misspecified fixed threshold. (2) \emph{Dynamic-Static Decomposition (DSD)}: dynamic areas are identified by the inability of static Gaussians to represent motion, then a buffer-based soft decomposition separates dynamic and static Gaussians, accelerating rendering and focusing deformation capacity on dynamic regions. It contributes the first event-inclusive 4D benchmark with synthetic and real-world dynamic scenes.

\textbf{Core point:} the event threshold is a learnable, jointly-optimized parameter, not a fixed constant---an early step toward differentiable event-camera physics.

\subsection{FastEventDGS (CVPR 2026)}
\label{sec:dynamic:fast}

FastEventDGS~\cite{dai2026fasteventdgs} is the first method to train Deformable 3DGS for dynamic scenes using \emph{only a single monocular event camera}, with no auxiliary RGB sensor. Its four design pillars are:
\begin{itemize}
    \item \textbf{Continuous trajectory parameterization:} cubic B-splines fit sparse poses into a continuous trajectory, enabling event supervision at any timestamp $\tau$---essential because event frequency vastly exceeds pose sampling rate.
    \item \textbf{Dual event generation model:} an Event Single Integral (ESI) loss provides global photometric consistency, while a linear event-flow loss $\Delta I(\mathbf{x})\approx-\nabla I\cdot\mathbf{v}\,\Delta\tau$ provides local geometric constraints via camera flow + Gaussian flow.
    \item \textbf{Local patch motion loss:} projects sampled Gaussian centers to the image plane and enforces that consecutive event integrals on patches around them are consistent, pinning Gaussian trajectories to event edges and suppressing background jitter.
    \item \textbf{Refinement stage:} uses the VGGT feed-forward model as a depth expert with a Scale-Invariant Logarithmic (SiLog) loss, plus non-event-region and spatial-smoothness regularizers.
\end{itemize}
It improves PSNR from $\sim$16\,dB to 22--24\,dB on synthetic and real fast-motion datasets.

\textbf{Core point:} events should be treated as explicit supervision for the motion field, not merely brightness differences; the combination of global integration + local flow + motion-edge pinning resolves the monocular under-constraint.

\subsection{ERF-GS (CVMJ 2026)}
\label{sec:dynamic:erfgs}

ERF-GS~\cite{bai2026erfgs} (the focus of this project's codebase) reconstructs fast motion from \emph{disjoint} multi-view Event-RGB viewpoints: some cameras supply only RGB, others only events, never both. Its three contributions are:
\begin{itemize}
    \item \textbf{linlog + STE event loss:} predicted events are the linlog-luma difference between two deformed-frame renders (Eq.~\eqref{eq:linlog}); the loss uses a straight-through estimator (STE) to quantize predictions to integer polarity counts while preserving gradients, normalized by $\max(|e_{gt}|,1)$.
    \item \textbf{Linear-motion + constant-color regularization:} between event-interval endpoints, the deformed position is required to follow a linear interpolation and the color to remain constant, encoding a constant-velocity prior.
    \item \textbf{Event-driven densification:} event pixels are back-projected at multiple candidate depths across views, filtered by multi-view consistency, selected by deformation-aware isolation, and initialized via motion-warped KNN. This spawns new Gaussians in fast-moving regions that RGB-gradient densification misses.
\end{itemize}
The disjoint-viewpoint setup forces cross-modal fusion and is the paper's defining experimental contribution.

\textbf{Core point:} events drive not only the loss but also the \emph{densification}---a capability unique to explicit GS representations.

\subsection{PEGS (arXiv 2025)}
\label{sec:dynamic:pegs}

PEGS~\cite{xu2025pegs} brings \emph{physics priors} to event-driven GS for large spatiotemporal-scale rigid-body motion. It layers three supervisions: (i) acceleration-consistency constraints from Newtonian dynamics, (ii) high-temporal-resolution event-flow guidance, and (iii) Kalman regularization fusing multi-source observations. A motion-aware simulated-annealing strategy schedules the constraints. It contributes the first paired RGB-Event dataset of natural fast rigid-body motion.

\textbf{Core point:} physics priors (acceleration consistency) and event supervision are complementary---events give high-rate positions, physics gives second-order smoothness---and their combination is, to our knowledge, unique to PEGS within the Event-GS family.

\subsection{Comparison of dynamic methods}

\begin{table}[t]
\centering
\caption{Comparison of dynamic/4D Event-GS methods.}
\label{tab:dynamic}
\small
\begin{tabular}{@{}p{2.3cm}p{1.1cm}p{1.1cm}p{1.6cm}p{1.0cm}@{}}
\toprule
Method & Input & Pose & Event model & Physics \\
\midrule
Event-Boosted~\cite{xu2025eventboosted} & E+R & Known & Learnable thresh. & -- \\
FastEventDGS~\cite{dai2026fasteventdgs} & E & Sparse & ESI+flow & -- \\
ERF-GS~\cite{bai2026erfgs} & E+R & Known & linlog+STE & Linear \\
PEGS~\cite{xu2025pegs} & E+R & Known & Event flow & Accel. \\
\bottomrule
\end{tabular}
\end{table}

Table~\ref{tab:dynamic} shows the diversity: only FastEventDGS is event-only; only PEGS uses physics; only ERF-GS uses event-driven densification; only Event-Boosted jointly models the threshold. This four-way design space---input modality, pose regime, event model, prior type---maps the open frontier of dynamic Event-GS.

\section{SLAM and Tracking}
\label{sec:slam}

Event-driven SLAM and tracking with Gaussian-Splatting maps is a smaller but growing category, sitting between the static-NVS and dynamic-4D lines.

\subsection{GS-EVT (ICRA 2025)}
\label{sec:slam:gsevt}

As discussed in \S\ref{sec:static:gsevt}, GS-EVT~\cite{liu2025gsevt} builds a 3DGS map from RGB frames and tracks an event camera against it, exploiting the complementary strengths of the two modalities: RGB for dense mapping, events for robust high-speed tracking.

\subsection{MotionGS-SLAM (arXiv 2026)}
\label{sec:slam:motiongs}

MotionGS-SLAM~\cite{hu2026motiongs} takes a conceptually distinctive approach: rather than treating motion blur as an inverse deblurring problem, it models blur \emph{forward} inside the renderer. Events modulate the Gaussian kernels in two ways: (i) \emph{spatial modulation} transforms isotropic points into anisotropic ellipses aligned with the motion direction, mimicking the blur trail; (ii) \emph{temporal modulation} adjusts exposure-integration sampling density by local velocity. A 4D spatiotemporal hash grid associates events with Gaussians. This yields a motion-blur-robust SLAM system that renders blur as a first-class phenomenon.

\textbf{Core point:} rather than remove blur, model it---events provide the per-pixel motion direction and magnitude that parameterize the blur kernel. This inverts the deblurring category's premise and connects SLAM to the forward-blur-modeling strand of EvFlow-GS and EBAD-Gaussian.

\section{Datasets and Simulation}
\label{sec:datasets}

A persistent bottleneck for Event-GS is data scarcity: real event-camera captures with paired RGB, poses, and geometry are expensive and geometrically inconsistent across views. One work directly addresses this.

\subsection{GS2E (NeurIPS 2025)}
\label{sec:datasets:gs2e}

GS2E~\cite{li2025gs2e} uses 3DGS itself as an event-data generator. It reconstructs real sparse multi-view RGB scenes with 3DGS, then synthesizes events via a physics-informed simulation pipeline with (i) adaptive trajectory interpolation for dense virtual viewpoints and (ii) physically consistent contrast-threshold modeling. This produces a large-scale, multi-view, geometrically consistent event dataset of 1150+ scenes, alleviating the data scarcity that limits event-based method development.

\textbf{Core point:} 3DGS is not only a consumer of event data but a \emph{producer} of it---closing a virtuous loop where GS reconstructions generate the training data for the next generation of event methods.

\subsection{Datasets contributed by method papers}

Beyond GS2E, several method papers contribute datasets that have become de facto benchmarks: EvaGaussians~\cite{yu2025evagaussians} (RGB+event+pose), Dark-EvGS~\cite{wu2025darkevgs} (real dark-light event-radiance), PEGS~\cite{xu2025pegs} (natural fast rigid-body RGB-Event), Event-Boosted~\cite{xu2025eventboosted} (event-inclusive 4D), and AsyncEvGS~\cite{dai2026asyncevgs} (asynchronous RGB-Event). The proliferation of these datasets is itself an indicator of the subfield's maturation.

\section{Cross-Cutting Technical Synthesis}
\label{sec:crosscutting}

Beyond the per-category narrative, the Event-GS literature converges on a small set of recurring technical dimensions. We extract them here to make the design space explicit.

\subsection{Event-to-Image vs.\ Event-to-Flow Supervision}
\label{sec:cross:supervision}

Two paradigms dominate event supervision:
\begin{itemize}
    \item \textbf{Event-to-image (intensity-domain):} render two frames, compute linlog-luma difference (Eq.~\eqref{eq:linlog}) or EDI integral (Eq.~\eqref{eq:edi}), compare to GT event voxel/integral. Used by E2GS~\cite{deguchi2024e2gs}, EvaGaussians~\cite{yu2025evagaussians}, EBAD-Gaussian~\cite{deng2025ebad}, ERF-GS~\cite{bai2026erfgs}, Dark-EvGS~\cite{wu2025darkevgs}. Mature and simple, but EvFlow-GS~\cite{an2026evflow} identifies four flaws: noise amplification, constant-threshold assumption, low-intensity degradation, and implicit constant-velocity bias.
    \item \textbf{Event-to-flow (motion-domain):} use the linear event generation model $\Delta I\approx-\nabla I\cdot\mathbf{v}\,\Delta\tau$ to supervise 2D optical flow or Gaussian screen motion, avoiding the intensity-domain altogether. Used by EvFlow-GS~\cite{an2026evflow} (optical flow warp), EF-3DGS~\cite{liao2025ef3dgs} (Contrast Maximization), FastEventDGS~\cite{dai2026fasteventdgs} (event-flow loss). Rising, because it sidesteps the noise and threshold issues.
\end{itemize}
The trend is toward flow-domain supervision or hybrid (global integration + local flow, as in FastEventDGS).

\subsection{Continuous-Time Trajectory Parameterization}
\label{sec:cross:trajectory}

Because event frequencies ($\sim$\,MHz) vastly exceed pose sampling rates (SfM/mocap), supervising only at discrete pose timestamps wastes event information. Three methods independently converge on \emph{continuous-time trajectory} parameterization: cubic B-splines (FastEventDGS~\cite{dai2026fasteventdgs}), B\'ezier SE(3) curves (EvFlow-GS~\cite{an2026evflow}, EBAD-Gaussian~\cite{deng2025ebad}), and continuous-time trajectory with temporally-coupled poses (EvTrajGS~\cite{chen2026evtrajgs}). This is becoming a standard component; any-time event supervision is now a prerequisite for fast-motion reconstruction.

\subsection{Deformation Fields}
\label{sec:cross:deformation}

Dynamic Event-GS methods inherit deformation backbones from RGB 4DGS: deformation MLPs with HexPlane~\cite{cao2023hexplane} features (ERF-GS~\cite{bai2026erfgs}, Event-Boosted~\cite{xu2025eventboosted}), or similar spatiotemporal grids (FastEventDGS~\cite{dai2026fasteventdgs}). A consistent design choice is to \emph{not} deform opacity and SH color by default (ERF-GS sets \texttt{no\_do=True}, \texttt{no\_no\_dshs=True}), only position/scale/rotation, to reduce under-constraint. Event-Boosted's dynamic-static decomposition is a complementary strategy: skip deformation entirely in static regions. The interplay between deformation capacity and event supervision is still under-explored---only four dynamic Event-GS methods exist as of mid-2026.

\subsection{Physical Priors}
\label{sec:cross:physics}

Physics constraints appear in two forms. PEGS~\cite{xu2025pegs} uses Newtonian acceleration consistency (second-order temporal smoothness) and Kalman fusion for rigid-body motion. ERF-GS~\cite{bai2026erfgs} encodes a constant-velocity (linear-motion) prior as a soft regularizer between event-interval endpoints. The gap: no method combines deformable (non-rigid) Gaussians with second-order physics priors, which would require differentiating the deformation field twice with respect to time---feasible but unexplored.

\subsection{Forward Blur Modeling}
\label{sec:cross:forwardblur}

A counter-current to the deblurring paradigm treats blur as a forward phenomenon to be \emph{modeled}, not removed. MotionGS-SLAM~\cite{hu2026motiongs} modulates Gaussian kernels by event-derived motion; EvFlow-GS~\cite{an2026evflow} and EBAD-Gaussian~\cite{deng2025ebad} synthesize blur from multiple sharp frames to construct a blur loss. This reframing---blur as observation, not noise---unifies the deblurring and SLAM strands.

\subsection{Event-Driven Densification}
\label{sec:cross:densify}

Unique to explicit GS representations, event-driven densification back-projects event pixels to spawn new Gaussians in fast-moving regions. ERF-GS~\cite{bai2026erfgs} is the only method to formalize this with multi-view consistency and motion-warped initialization. This is a capability with no analog in NeRF-based event methods and a likely growth area.

\subsection{Threshold and Sensor Modeling}
\label{sec:cross:threshold}

The event contrast threshold $C$ is in reality pixel-, time-, and scene-dependent, yet most methods fix it (ERF-GS uses $C=0.2$). Event-Boosted~\cite{xu2025eventboosted} jointly models the threshold with the GS field, and EvFlow-GS~\cite{an2026evflow} replaces the fixed-threshold model entirely. A fully differentiable event-camera physics module---learnable per-pixel threshold, noise, and response---remains open (see \S\ref{sec:discussion}).

\section{Discussion and Open Problems}
\label{sec:discussion}

The Event-GS literature has matured rapidly but several gaps remain conspicuously open. We identify five, ordered roughly by modality-specificity (most event-unique first).

\subsection{Occlusion Handling via Event Appearance/Disappearance Boundaries}
\label{sec:disc:occlusion}

No Event-GS method explicitly handles occlusion. Existing dynamic methods (FastEventDGS~\cite{dai2026fasteventdgs}, ERF-GS~\cite{bai2026erfgs}, Event-Boosted~\cite{xu2025eventboosted}) assume Gaussians remain visible; occluded Gaussians drift or are mis-pruned for lack of supervision. Events offer a modality-unique signal: a pixel stops triggering events when its object is occluded and bursts when it reappears. The ``appearance/disappearance boundaries'' of event activity could online-detect occlusion intervals, during which occluded Gaussians are frozen (no deformation, no densification), resuming optimization on reappearance. RGB cannot distinguish ``object disappeared'' from ``object occluded''; events can, because a statically occluded object triggers no events. This is the most event-unique open direction.

\subsection{Static-Prior Bootstrapping for 4DGS}
\label{sec:disc:bootstrap}

All dynamic Event-GS methods cold-start from sparse COLMAP point clouds, learning appearance, motion, and pose simultaneously. None exploits the now-availability of high-quality static 3DGS reconstructions as a prior. A ``bootstrap'' paradigm---initialize \texttt{GaussianModel} from a converged static 3DGS \texttt{.ply} (xyz+SH+scale+opacity+rotation), zero-initialize the deformation field to identity, and let events supervise only the motion increment---would lock appearance/geometry and free optimization capacity for motion. This directly addresses the cold-start failure mode where fast-moving objects ``disappear'' because static-region Gaussians outcompete them for gradient.

\subsection{Per-Event Point-Process Losses}
\label{sec:disc:pointprocess}

Every Event-GS method \emph{aggregates} events into voxels, frames, or integrals, discarding microsecond-level time resolution. Worse, GT aggregation counts total windowed activity while prediction examines only endpoint net change, systematically mismatching oscillatory motion. The event stream is a marked temporal point process: each $(x,y,t,p)$ is a physical constraint ``log-luma crossed $\pm C$ at time $t$.'' A per-event loss---rendering at exact event timestamps and enforcing the threshold-crossing constraint---would be the most physically faithful supervision form. No Event-GS method implements this; it requires preserving event timestamps (current preprocessing discards them), per-event sampling, and a differentiable point-process likelihood. The academic novelty is high; the engineering cost is moderate.

\subsection{Differentiable Event-Camera Physics Simulator}
\label{sec:disc:sim}

The linlog + fixed-threshold + STE pipeline (ERF-GS~\cite{bai2026erfgs}) is a hard-quantization approximation with poor gradient fidelity. GS2E~\cite{li2025gs2e} uses ESIM for simulation but it is not differentiable or jointly trained. A fully differentiable event-physics module---sigmoid-soft threshold crossing, learnable per-pixel threshold/bias/noise, jointly trained with 3DGS---would let gradients flow through the true event-generation physics rather than a straight-through surrogate. This is a general-purpose enhancement stackable onto any of the above.

\subsection{Physics-Constrained Deformable Event-GS}
\label{sec:disc:physics}

PEGS~\cite{xu2025pegs} demonstrates physics priors for rigid SE(3) motion but does not deform Gaussians; ERF-GS/FastEventDGS deform Gaussians but lack physics. Combining a deformable event-GS with (i) acceleration-consistency regularization (second derivative of the deformation field w.r.t.\ time, feasible since deformation fields are continuously differentiable) and (ii) local-rigidity priors (KNN-neighbor relative-displacement L2) would constrain the under-determined motion field. The challenge is weighting physics vs.\ event loss; the benefit is temporal smoothness that events alone cannot enforce.

\subsection{Evaluation and Benchmarks}
\label{sec:disc:benchmarks}

The field lacks a unified benchmark. Each method paper contributes its own dataset (EvaGaussians, Dark-EvGS, PEGS, Event-Boosted, AsyncEvGS) with inconsistent metrics, masks, and splits. A community benchmark---with standardized dynamic/static region metrics (cf.\ ERF-GS's Dynamic-* / Static-* metrics using dynamic masks), shared test scenes across deblurring/4D/SLAM, and a common event-simulation protocol---would accelerate progress. GS2E's 1150+ scenes are a starting point for synthetic evaluation; real-data standardization remains open.

\section{Conclusion}
\label{sec:conclusion}

Event cameras and Gaussian Splatting are a natural pairing: the explicit, real-time, differentiable rasterizer of 3DGS is an ideal substrate for the high-temporal-resolution, motion-rich, but intensity-poor signal of events. In under three years, the Event-GS subfield has grown from a single proof-of-concept~\cite{han2024event3dgs,deguchi2024e2gs} into a multi-track research program spanning static NVS, deblurring, low-light/HDR, dynamic 4D, SLAM, and simulation, with dedicated venues at CVPR, ICCV, ECCV, NeurIPS, and ICML.

This survey organized $24+$ methods into a six-category taxonomy and analyzed the core contribution, event-generation model, and supervision form of each. The cross-cutting synthesis reveals convergence on continuous-time trajectory parameterization, a shift from intensity-domain to flow-domain event supervision, and the emergence of event-driven densification as a GS-unique capability. The most promising open directions---occlusion-aware 4DGS via event boundaries, static-prior bootstrapping, per-event point-process losses, differentiable event-camera physics, and physics-constrained deformation---are precisely those that exploit the event camera's modality-unique properties beyond what RGB can provide.

As event-camera hardware matures and 3DGS becomes the default real-time radiance-field representation, we expect Event-GS to become a standard component of future AR/VR, robotics, and high-speed imaging systems rather than a niche subfield. The next wave of progress will likely come from replacing the linlog+STE approximation with differentiable event physics, and from treating events as explicit motion-field supervision rather than aggregated intensity differences.


\bibliographystyle{IEEEtran}
\bibliography{references}

\end{document}
